\section{방법론}
\label{sec:method}

\rev{이 장은 그림~\ref{fig:pipeline}의 순서대로 표본 구성(A, B), 세 연구 질문의 모델(C), 비교 방법(D)을 설명한다.}

\begin{figure*}[t]
\centering
\includegraphics[width=\textwidth]{fig_pipeline_clean.pdf}
\caption{데이터 파이프라인과 모델링·평가 절차}
\label{fig:pipeline}
\end{figure*}

%======================================================================
\subsection{데이터 수집}
\label{sec:method:data}
PUBG의 공개 텔레메트리는 경기에 참가한 모든 플레이어의 위치를 약 10초마다 기록하고, 피해, 기절, 처치, 아이템, 안전구역의 상태, 단계 전환 같은 사건을 일어난 시각과 함께 기록한다~\cite{PUBGAPI}.
본 연구는 공개 API가 제공하는 직전 24시간의 표본 경기 목록을 2026년 3월 27일--6월 2일(UTC)에 35회 조회하여, 3월 25일--6월 1일에 생성된 스팀 서버의 경기를 수집하였다.
이 기간의 게임 버전은 PC 업데이트 40.2--41.2이며~\cite{PUBGPatchNotes2026}, 일부 맵 요소를 빼면 안전구역 규칙은 같고 규칙 개편~\cite{PUBGBlueZone2026}은 수집 이후에 적용되었다.
\rev{이처럼 촘촘한 경기 내 위치 기록을 담은 공개 자료는 드물고 기존 연구가 가공한 자료도 공개되지 않았다~\cite{Sagong25a,KRAFTON}. 본 연구의 코드는 게재 후 공개한다.
분석에는 공개 API의 경기 기록만 썼고 개인을 식별하는 분석은 하지 않았으며, 경기와 표본의 규모는 표~\ref{tab:cohorts}와 같다(구성한 팀--단계 표본은 2.9\,GB).}

%======================================================================
\subsection{\rev{예측 시점 기준 표본 구성}}
\label{sec:method:samples}
예측 시점보다 뒤에 일어난 일이 입력에 조금이라도 섞이면, 모델은 결과를 미리 보고 맞히는 셈이 된다.
\rev{예를 들어(그림~\ref{fig:pipeline}의 B 가운데 시간축) 어느 경기에서 3번째 구역으로의 축소가 1,200초에 끝나고 4번째 구역이 공개되면, 단계 3의 예측 시점은 1,200초, 특징 절단 시각은 1,190초이다.
팀 A의 입력은 1,190초 이전의 팀원 위치와 생존 여부, 주변 팀과의 거리와 교전, 구역에 대한 위치로만 만들며, 1,195초의 교전은 쓰지 않는다.
팀 A가 1,200초에 살아 있으면 단계 3의 표본이 되고, 4번째 구역으로의 축소가 끝나기 전에 전멸하면 라벨이 1, 아니면 0이다.}

\subsubsection{시간축, 예측 시점, 입력 절단}
\rev{모든 기록의 시각은 경기 시작부터 흐른 초로 통일한다.}
경기 $m$의 단계 $k$는 $k$번째 안전구역이 공개된 때부터 그 구역으로의 축소가 끝날 때까지이다.
\rev{축소가 끝나 다음 안전구역이 공개되는 순간을 예측 시점 $b_{m,k}$로 두고,} 입력은 그보다 $\Delta=10$초 앞선 특징 절단 시각 $a_{m,k}=b_{m,k}-\Delta$ 이전의 기록만으로 만든다.
\rev{$\Delta$는 약 10초인 위치 기록 간격과 위치 보간의 여유를 고려한 값이며,} 이 규칙은 생존 판정, 팀 간 관계, attention 마스크에도 똑같이 적용된다.
처치된 플레이어는 \rev{처치 20초 이후 $a_{m,k}$ 이전에} 체력이 남은 위치 기록이 다시 나타날 때만 살아 돌아온 것으로 본다.
\rev{코호트마다 17경기, 모두 204경기의 1,293개 경계에서 $a_{m,k}$ 이후의 기록을 지우고 입력을 다시 만들어도 원래 입력과 같았고, 절단을 뺀 빌더는 별도 누수 시험의 21개 경계에서 모두 검출되었다.}

\subsubsection{사건 기반 라벨과 예측 대상}
정답, 즉 라벨은 “다음 단계가 끝나기 전에 그 팀이 전멸하는가”이다.
기절한 플레이어는 다시 일어날 수 있으므로 전멸로 보지 않으며, 경기 $m$에서 팀 $i$의 탈락 시각 $T_{m,i}$는 다시 돌아오지 않은 마지막 팀원이 처치된 시각이다(우승 팀은 $T_{m,i}=\infty$, 경기 종료 시각 $C_m$까지 관측).
단계 $k$에서 예측의 대상은 예측 시점에 아직 살아 있는 팀들, 즉 위험집합 $\mathcal{R}_{m,k}=\{i: T_{m,i}>b_{m,k}\}$이며, \rev{따라서 절단 시각과 예측 시점 사이 10초 안에 탈락한 팀은 그 단계의 표본에서 빠진다.} \rev{$T_{m,i}$는 이후 기록을 쓰므로 위험집합을 절단 시각 이전 기록만으로 판정한 생존 상태와 대조하였고, 817,766개 표본 가운데 살아 있는 팀원이 없던 표본은 2개였으며 시험 표본에는 없었다.}
각 팀의 라벨은 $Y_{m,i,k}=\mathbb{I}[T_{m,i}\le b_{m,k+1}]$이며, $\mathbb{I}[\cdot]$은 괄호 안이 참이면 1, 거짓이면 0이다. 단계 $k+1$이 경기의 마지막 단계이면 $b_{m,k+1}=C_m$으로 둔다.
표본 $(m,i,k)$는 단계 $k=1,\dots,7$의 경계에서 만들며(솔로 경기에서는 팀이 곧 플레이어), 모델이 추정할 양은 다음 단계 탈락 위험
\begin{equation}
\lambda_{m,i,k}=\Pr\!\left(Y_{m,i,k}=1 \mid i\in\mathcal{R}_{m,k},\,\mathcal{H}_{m}(a_{m,k})\right)
\label{eq:hazard}
\end{equation}
이다. 여기서 $\mathcal{H}_{m}(t)$는 경기 $m$에서 시각 $t$ 이전에 기록된 텔레메트리 전체이다.
이는 기준 시점의 생존자만 놓고 다음 구간의 사건확률을 보는 landmark 방식의 이산시간 hazard이다~\cite{SingerWillett93,vanHouwelingen07}.

\subsubsection{코호트와 분할}
\label{sec:method:cohort}
맵, 시점, 팀 규모가 다르면 이동 거리와 교전 양상이 달라지므로, 같은 조합의 경기끼리 묶어 코호트로 분석한다.
분석 대상은 일반 경기와 랭크 경기이며, 튜토리얼·훈련장 경기, 계정 대부분이 AI인 캐주얼 모드, 비플레이어 전투 계정이 많은 Vikendi 맵은 실제 경쟁과 달라 제외하였다.
일반 경기에도 AI 계정이 섞여 있으므로 각 팀의 AI 계정 비율을 입력에 넣었다.
경기가 충분한 12개 조합을 코호트로 사용하였고(표~\ref{tab:cohorts}), 같은 경기의 팀들이 학습과 시험에 나뉘어 들어가지 않도록 경기 단위로 학습, 검증, 시험 집합을 \rev{약 8:1:1로} 고정 분할하였다.

\begin{table}[t]
\centering
\caption{코호트별 경기 수, 팀--단계 표본 수, 다음 단계 탈락률}
\label{tab:cohorts}
\footnotesize
\setlength{\tabcolsep}{2.7pt}
\begin{tabular}{lccrrrr}
\hline
맵 & 시점 & 인원 & 경기 & 학습/검증/시험 & 표본 & 탈락률(\%) \\
\hline
Erangel & 3인칭 & 1 & 224 & 189/15/20 & 31,740 & 35.1 \\
Erangel & 3인칭 & 2 & 656 & 526/65/65 & 59,170 & 29.5 \\
Taego & 3인칭 & 2 & 573 & 459/57/57 & 47,223 & 29.2 \\
Erangel & 1인칭 & 4 & 430 & 344/43/43 & 21,849 & 27.8 \\
Erangel & 3인칭 & 4 & 3,283 & 2,638/321/324 & 196,442 & 25.7 \\
Karakin & 3인칭 & 4 & 333 & 267/33/33 & 12,669 & 32.1 \\
Miramar & 3인칭 & 4 & 1,543 & 1,235/154/154 & 92,279 & 24.6 \\
Paramo & 3인칭 & 4 & 333 & 267/33/33 & 9,963 & 30.5 \\
Rondo & 3인칭 & 4 & 1,039 & 833/103/103 & 53,683 & 25.2 \\
Sanhok & 3인칭 & 4 & 940 & 752/94/94 & 48,793 & 27.6 \\
Taego & 1인칭 & 4 & 435 & 349/43/43 & 23,390 & 30.0 \\
Taego & 3인칭 & 4 & 4,159 & 3,329/415/415 & 220,565 & 24.9 \\
\hline
합계 & & & 13,948 & 11,188/1,376/1,384 & 817,766 & 26.6 \\
\hline
\end{tabular}
\\[2pt]
\parbox{\columnwidth}{\scriptsize 2026년 3월 25일--6월 1일(UTC) 생성 경기(업데이트 40.2/41.1/41.2: 4,439/6,971/2,538경기). 텔레메트리 기록 약 3억 2,240만 건(위치 9,643만, 사건 2억 2,597만).}
\end{table}

%======================================================================
\subsection{설계와 분석}
\label{sec:method:design}
\rev{서식지 맥락 표현을 먼저 정의한 뒤 세 분석을 연구 질문 순서로 설명한다(그림~\ref{fig:pipeline}의 C).}

\subsubsection{서식지 맥락 표현}
\label{sec:method:representation}
줄어드는 서식지에서 생물 집단의 존속이 위치와 경쟁에 달려 있듯이, 팀의 생존도 구역에 대한 위치와 주변의 붐빔에 달려 있다고 보고 팀의 상태를 네 가지로 나타낸다.
구역 경계에 대한 노출은 팀이 \rev{현재 구역(블루존)} 밖으로 얼마나 나가 있는지(구역 밖 깊이) 또는 안쪽으로 얼마나 들어와 있는지(구역 안쪽 깊이)이다.
이동 부담은 \rev{절단 시각에 보이는 목표 안전구역에 들어가는 데} 필요한 속도가 팀이 실제로 움직인 속도보다 얼마나 큰지\rev{(이동 부족분)}이며~\cite{Nathan08}, \rev{예측 시점에 새로 공개되는 구역은 입력에 쓰지 않는다.}
주변 경쟁 밀도는 팀 주변 200\,m 안에 있는 다른 팀의 수이고, 남은 팀원 수는 절단 시각에 살아 있다고 판정된 팀원의 수이다.
\rev{본 연구의 입력은 이 구성개념을 담은 서식지 맥락 블록과 나머지 팀 특징 블록으로 이루어지며(표~\ref{tab:blocks}),} 서식지 맥락 블록은 RQ1의 해석 모델과 분해 단계 (a1)에, 나머지 팀 특징은 (a2)에, 관계 그래프는 PTR-Former에 쓰인다.
또한 같은 단계의 팀들 사이에는 가까이 있는 관계, 서로 싸운 관계, 함께 구역의 압박을 받는 관계의 세 가지를 정의한다.

\begin{table}[t]
\centering
\caption{입력 블록과 대표 특징}
\label{tab:blocks}
\footnotesize
\setlength{\tabcolsep}{3pt}
\renewcommand{\arraystretch}{1.0}
\begin{tabular}{p{1.8cm}p{5.4cm}r}
\hline
블록 & 대표 특징 & 수 \\
\hline
\multicolumn{3}{l}{\emph{본 연구}} \\
\textbf{서식지 맥락} & \textbf{구역 밖·안쪽 깊이, 경계까지 거리, 필요 속도와 이동 부족분, 반경 100--500\,m의 적 팀 수, 구역 크기와 축소율, 진입 혼잡도} & 44 \\
경기 맥락 & 단계 번호, 단계 길이, 팀 규모, 생존 팀 수 & 6 \\
팀 상태 & 생존 인원, 평균·최소 체력, 팀원 간 거리 & 9 \\
이동 & 경로 길이, 평균 속도, 굴곡도, 구역 진입 시점 & 7 \\
교전 & 가한·받은 피해, 처치·기절, 교전 상대 수 & 13 \\
자원 & 아이템 획득·장착·사용, 회복, 차량 이용 & 16 \\
공간·고도 & 가장 가까운 적까지 거리, 적 밀도, 고도 & 12 \\
관계 요약 & 근접·교전·공동 압박 관계의 연결 수와 강도 & 43 \\
관계 그래프 & 근접, 교전(보낸·받은), 공동 압박 가중치 & --- \\
\hline
\multicolumn{3}{l}{\emph{기존 설계}} \\
Sagong 등 & 단계별 공격·처치·피해·어시스트 수, 이동 거리, 밀집도, 아이템, \rev{보급 상자,} 교전·활동 점수, 플레이 유형, 유형별 생존율 & 12 \\
KRAFTON & 처치·기절, 무기별 피해, 받은 피해와 회복, 장비 수준, 주변 아군·적 수, 구역까지 거리, 팀 체력과 생존 인원, 팀원 위치 & 55+ \\
\hline
\end{tabular}
\end{table}

\subsubsection{RQ1: 생태 hazard 모델}
\label{sec:method:rq1}
RQ1은 각 요인이 위험을 어느 방향으로 얼마나 바꾸는지를 묻기 때문에, 계수를 그대로 읽을 수 있는 로지스틱 회귀
\begin{equation}
\mathrm{logit}\,\lambda_{m,i,k}=\alpha_k+\textstyle\sum_{c\in\mathcal{C}}\beta_c z_{c,m,i,k}
\label{eq:ehm}
\end{equation}
를 코호트마다 추정한다.
여기서 $\mathrm{logit}\,p=\log\{p/(1-p)\}$이고, $\alpha_k$는 단계 $k$의 \rev{기저위험, 곧 그 단계의 모든 팀에 공통인 위험 수준}(절편)이다.
$\mathcal{C}$는 구역 밖 깊이, 구역 안쪽 깊이, 이동 부족분, 주변 경쟁 밀도, 남은 팀원 수의 다섯 변수와 통제 변수인 절단 시각의 생존 팀 수의 로그로 이루어지며, 솔로 코호트에서는 남은 팀원 수를 뺀다.
$z_{c,m,i,k}$는 변수 $c$를 그 코호트의 학습 집합 평균과 표준편차로 표준화한 값이고, $\beta_c$는 그 변수가 한 표준편차 늘 때 위험의 로짓이 바뀌는 양이다.
표준오차는 같은 경기의 표본끼리 묶어 계산하고, 12개 코호트의 계수는 확률효과 모형으로 합친다~\cite{DerSimonian86}.
\rev{단계에 따른 변화는 단계 구간(초반 1--2, 중반 3--5, 후반 6--7)마다 구간 안 표준편차 단위의 계수를 추정하고 구간 간 동일성을 다변량 확률효과 Wald 검정으로 보며, \rev{공통 단위(코호트 학습 집합의 표준편차)로도 같은 검정을 한다.}
또한 단계 절편과 네 변수(필요 속도, 경쟁 밀도, 생존 팀 수의 로그, 남은 팀원 수)만 쓴 간이 모형이 GBDT 판별력의 얼마를 재현하는지를 $(\mathrm{AUC}-0.5)$의 비로 잰다.} \rev{이 간이 모형은 탐색적 분석이며, 바깥 깊이가 필요 속도와 거의 공선이어서($r=0.89$--$0.95$) 두 깊이 항을 빼고 생존 팀 수의 로그는 경기 수준 통제 변수로 둔다.}

\subsubsection{RQ2: 본 연구 설계와 기존 설계}
\label{sec:method:rq2}
본 연구의 설계는 \rev{단계를 조건으로} “다음 단계에 탈락할 확률” 자체를 맞히도록 학습하며, 위험을
\begin{equation}
\mathrm{logit}\,\lambda_{m,i,k}=\alpha_k+f(\mathbf{x}_{m,i,k},\mathcal{G}_{m,k})
\label{eq:logit}
\end{equation}
로 둔다. 여기서 $\mathbf{x}_{m,i,k}$는 팀 $i$의 서식지 맥락을 포함한 입력 특징 벡터, $\mathcal{G}_{m,k}$는 $a_{m,k}$에 살아 있는 팀들 사이의 세 관계를 담은 그래프, $f$는 학습되는 함수이다.
모델은 학습 집합의 모든 팀--단계 표본에 대한 교차엔트로피 $-\sum[Y_{m,i,k}\log\lambda_{m,i,k}+(1-Y_{m,i,k})\log(1-\lambda_{m,i,k})]$를 최소화한다.
\rev{식~\eqref{eq:logit}은 개념적 형태이며, 두 학습기는 $\alpha_k$를 단계 정보 입력으로 표현하고 별도 절편 모수는 분해의 단계 (d)에서만 더한다.}
하나는 단계 번호를 입력으로 받는 그래디언트 부스팅 결정나무(GBDT)~\cite{Friedman01}이고, 다른 하나는 구조의 기여를 분해하기 위한 참조 모델인 PTR-Former\rev{(phase--team--relation Transformer)}이다.
PTR-Former는 같은 팀의 이전 단계들을 돌아보는 attention(단계 축)과 같은 단계의 다른 생존 팀들을 함께 보는 attention(팀 축)을 번갈아 쓰고~\cite{Vaswani17}, 팀 축 attention에는 관계의 종류에 따라 학습되는 편향을 더하며~\cite{Ying21}, \rev{순위 보조 손실(가중치 0.1)을 함께 학습한다}\footnote{GBDT는 결정나무 200개와 학습률 0.06을, PTR-Former는 차원 128, attention 헤드 4개, 블록 2개, dropout 0.1, AdamW(학습률 $10^{-3}$)와 코사인 학습률 감소를 사용하였고, 최대 12 epoch 동안 검증 로그 손실을 기준으로 조기 종료하였다. \rev{모든 코호트에서 같은 초기화 하나로 학습하였다.}}.

KRAFTON의 예측기는 이와 달리 “누가 먼저 탈락하는가”라는 순서를 배운다.
구체적으로는 Cox 부분우도의 항 $\exp(s_i)/\sum_{j\in\mathcal{R}}\exp(s_j)$에 가중치를 둔 목적을 학습하며~\cite{Cox72}, $s_i$는 한 시점에서 팀 $i$의 점수, $\mathcal{R}$은 그 시점에 살아 있는 팀의 집합이다.
모든 팀의 점수에 같은 값을 더해도 이 항은 그대로이므로, 이 목적만으로는 “단계 4가 단계 2보다 전체적으로 더 위험하다”는 정보, 곧 식~\eqref{eq:logit}의 $\alpha_k$가 \rev{식별되지 않는다}.
그래서 점수에 단계별 \rev{기저위험}을 덧붙인 확률 $1-\exp\{-\exp(\gamma_k+s_{m,i,k})\}$도 함께 비교한다. 여기서 $s_{m,i,k}$는 예측 시점 $b_{m,k}$에서 팀 $i$의 점수이고, $\gamma_k$는 \rev{검증 집합}의 위험집합에서 관측 탈락과 점수로 추정한 단계 $k$의 \rev{기저위험}이다~\cite{PrenticeGloeckler78}.

두 기존 설계는 원래의 입력, 학습 목적, 학습 방식을 그대로 두고 같은 분할과 같은 위험집합에서 평가하였으며, 차이와 적용 방식은 표~\ref{tab:compare}에 정리하였다.
또한 적용판을 우리 과제에 불리하게 만들지 않았는지 확인하기 위해, 각 적용판을 원래 과제의 지표(우승 팀 정확도와 탈락 순서 C-index, 솔로 코호트의 개인 생존 분류)로도 평가한다.

\begin{table*}[t]
\centering
\caption{기존 설계와 본 연구 설계의 비교}
\label{tab:compare}
\footnotesize
\renewcommand{\arraystretch}{1.1}
\begin{tabular}{p{2.3cm}p{4.3cm}p{4.6cm}p{4.6cm}}
\hline
 & Sagong 등~\cite{Sagong25a} & KRAFTON~\cite{KRAFTON} & 본 연구 \\
\hline
분석 단위 & 플레이어 & 팀 & 팀(솔로 경기에서는 개인) \\
예측 대상 & 다음 반 단계의 생존 여부 & 탈락 순서의 상대위험(우승확률로 변환) & \textbf{다음 단계 안에 탈락할 조건부 확률} \\
예측 시점 & 반 단계 경계 & 경기 중 여러 시점 & 단계 경계(다음 안전구역 공개 시점) \\
단계별 위험 수준 & 플레이 유형별 전이 생존율 & 표현하지 않음 & \rev{\textbf{단계 조건부 위험 수준}(단계를 입력으로)} \\
팀 간 상호작용 & 다루지 않음 & 팀 간 attention & \rev{참조 모델에만:} 팀 축·단계 축 attention과 관계 편향 \\
학습 목적 & 생존 분류 & 가중 Cox 부분우도 & \textbf{다음 단계 탈락의 교차엔트로피} \\
\hline
본 연구의 적용 & 2·4인 팀은 팀 단위로 집계, 예측 시점을 단계 경계로 통일 & 입력을 본 연구의 텔레메트리로 재구성\rev{(일부 근사)}, 출력을 \rev{검증 집합에서} 확률로 보정 & --- \\
\hline
\end{tabular}
\end{table*}

\subsubsection{RQ3: \rev{단계별 분해}}
\label{sec:method:decomposition}
RQ3는 성능 차이가 어디서 오는지를 묻는다. \rev{그래서 설계를 한 단계씩 바꾸며 성능 변화를 비교하고, 각 단계의 변화를 앞 단계 설계를 조건으로 한 그 단계의 기여로 읽는다. 대부분의 단계는 한 요소만 바꾸지만, KRAFTON 경로의 첫 단계는 입력과 학습 행을, (e)는 구조와 학습 설정을, (f)는 단계 축과 관계 편향을 함께 바꾼다.}
입력의 효과는 학습기를 GBDT로 고정하고 측정한다. Sagong 등의 입력에 서식지 맥락을 더하고(a1), 나머지 입력을 본 연구의 입력으로 바꾸며(a2), 따로 KRAFTON 입력을 본 연구의 입력으로 바꾼다(b).
학습 목적과 \rev{기저위험}의 효과는 KRAFTON의 구조에 본 연구의 입력을 넣고 Cox 목적으로 학습한 모델에서 출발하여, 단계별 \rev{기저위험}을 덧붙이고(c1), 학습 목적을 다음 단계 교차엔트로피로 바꾸고(c2), 단계 절편을 명시적으로 더하며(d) 측정한다.
구조의 효과는 KRAFTON 구조를 팀 축만 남긴 PTR-Former로 바꾸고(e), 단계 축과 관계 편향을 더하여(f) 측정하며, 두 단계에서는 순위 보조 손실을 쓰지 않는다.
마지막으로 GBDT에 KRAFTON 입력을 더하여 두 표현이 서로 보완하는지 확인한다(g).
\rev{이와 별도로 단계 축과 팀 축을 켜고 끄는 2$\times$2 실험(두 축의 AUC와 팀 축의 순위 상관), 관계를 무작위로 섞은 대조 모델, 서식지·구역 특징을 뺀 조건에서 구역을 별도 토큰으로 넣은 모델의 다섯 비교를 실험 전에 정해 두었다.
각 단계의 변화는 통합 AUC와 함께 단계 안·경기--단계 내부 AUC로도 잰다.}

%======================================================================
\subsection{평가}
\label{sec:method:evaluation}
모든 설계는 같은 시험 표본에서 같은 지표로 평가하고, 그림~\ref{fig:pipeline}의 D와 같이 코호트별 차이를 하나로 합친다.

\subsubsection{평가 지표}
주 지표인 AUC는 실제로 탈락한 팀에 탈락하지 않은 팀보다 높은 위험을 줄 확률이다.
단계가 뒤로 갈수록 모든 팀의 위험이 함께 올라 단계만 맞혀도 AUC가 높아질 수 있으므로, 같은 단계끼리만 비교한 단계 안 AUC\rev{(단계별 AUC의 평균)}와 같은 경기·같은 순간의 생존 팀끼리만 비교한 경기--단계 내부 AUC\rev{(두 부류가 함께 있는 위험집합의 쌍을 모아 계산)}를 함께 본다.
예측확률의 정확도는 로그 손실과 Brier 점수로 평가한다.
단계마다 달라지는 위험 수준을 얼마나 잘 나타내는지는 단계별 보정 오차 $\sum_{k}(n_k/n)\,|\bar{p}_k-\bar{y}_k|$\rev{를 코호트별로 구해 평균하여} 평가한다. 여기서 $n_k$는 단계 $k$의 시험 표본 수, $n=\sum_k n_k$, $\bar{p}_k$와 $\bar{y}_k$는 그 표본의 평균 예측확률과 실제 탈락률이다. \rev{단계 안의 보정은 각 단계의 시험 표본을 예측확률 $[0,1]$의 등간격 10개 구간 $j$로 나누어 구간별로 같은 양을 구한 기대 보정 오차(ECE) $\sum_{k}\sum_{j}(n_{kj}/n)\,|\bar{p}_{kj}-\bar{y}_{kj}|$로 보며, 이 역시 코호트별로 구해 평균한다.}

\subsubsection{통계적 비교}
두 설계의 차이가 우연인지 보기 위해, 시험 경기를 2,000번 다시 뽑는 경기 단위 쌍대 bootstrap으로 코호트마다 신뢰구간을 구한다.
12개 코호트의 결과는 코호트마다 효과가 다를 수 있음을 허용하는 DerSimonian--Laird 확률효과 모형으로 합치고~\cite{DerSimonian86}, REML 추정과 Hartung--Knapp 구간으로 같은 결론이 나오는지 확인한다~\cite{HartungKnapp01}.
\rev{실험 전에 정한 기준에 따라, 차이가 단측 검정에서 $p<0.05$이고 추정치가 0.005(순위 상관은 0.01) 이상이면 개선, 90\% 신뢰구간이 $\pm0.005$(순위 상관은 $\pm0.01$) 안에 있으면 실질적 동등, 그 밖은 불확정으로 판정한다.}
미리 정한 다섯 비교에는 Holm 보정을 적용하며~\cite{Holm79}, 나머지 비교는 탐색적 분석으로 보고한다. \rev{판정 규칙, SESOI, 다섯 비교는 본 실험 전에 날짜가 남는 설계 커밋으로 고정하였다. SESOI 0.005는 순서를 맞힌 쌍의 비율로 0.5\%p이며, 같은 자료의 이전 탐색 분석에서 본 작은 관계 효과와 그보다 큰 서식지 효과 사이의 값이다.}
끝으로 가장 큰 두 코호트에서 학습 초기화를 바꾸고 \rev{최대 학습 횟수를 두 배로 늘려} 결과가 유지되는지 확인한다.
